\documentclass[10pt,a4paper]{amsart}
\usepackage[margin=19mm]{geometry}
\usepackage{amsmath,amssymb,mathtools}
\usepackage[scr=rsfs,cal=euler]{mathalfa}
\usepackage{xfrac}
\usepackage[unicode,hidelinks,bookmarks=false]{hyperref}
\newcommand{\ie}{i.e.\ }
\newcommand{\cf}{cf.\ }
\newcommand{\resp}{respectively\ }
\newcommand{\Subsection}{\subsection}
\providecommand{\nobreakdash}{\nobreak\hbox{-}}
\setlength{\emergencystretch}{2em}
\multlinegap0pt
\usepackage{amssymb}
\usepackage{mathtools}
\usepackage{tikz}
\newtheorem{theorem}{Theorem}
\newtheorem{lemma}{Lemma}
\newcommand{\Ke}{K_7^{=}}
\newcommand{\dd}{\partial}
\title{A sharp density bound for 5-connected graphs\\with no $\Ke$ minor}
\author{Caibing Chang, Zijian Deng, Qinfei Tang, Caihong Yang}
\date{Source: arXiv:2609.26041v1 (CC BY 4.0)}
\begin{document}
\maketitle

\noindent\emph{Source and licence.} A sharp density bound for 5-connected graphs\\with no $\Ke$ minor. Version arXiv:2609.26041v1 (CC BY 4.0), distributed by arXiv under the Creative Commons Attribution 4.0 International licence, http://creativecommons.org/licenses/by/4.0/, as recorded in the arXiv licence metadata for this exact version and shown on the abstract page https://arxiv.org/abs/2609.26041v1. Full licence notice: \texttt{licenses/arxiv-2609.26041-CC-BY-4.0.txt}.\par\smallskip

\noindent\textit{Editorial scope.} Counted unit: the article's lemma lem:criticalirreducible (source0.tex lines 765--767). The article's complete original proof of it is delivered with it (source0.tex lines 768--795, one \texttt{proof} environment of 28 lines). No mathematics is written, repaired or re-derived: the statement and the proof are the article's own lines, in the article's own notation, and no retained line is substituted or re-flowed.  Retained uncounted as the source-local context the proof needs: lines 90--92; lines 217--222; lines 760--762.  Nothing was omitted from any retained span, so the package carries no omission record.  No retained line contains a citation, so the wrapper carries no bibliography and no bibliography entry.  No retained line contains a figure environment, an image-inclusion command or drawing code, so no image is delivered and the package declares no asset dependency.  Editorial formatting only, in addition: an AMS \texttt{amsart} preamble that restates the article's own declarations and macros used by the retained lines, namely the environments \texttt{theorem}, \texttt{lemma} on separate counters, together with the standard packages those lines need.  The retained environments print in this document as Theorem 1, Lemma 1 in order of appearance, whereas the article prints the same items with the numbering of its own full document.  The complete native source of the article as distributed in the arXiv e-print for this exact version is bundled unchanged as \texttt{sources/211341/source0.tex}.
\begin{theorem}\label{thm:aux}
Every $4$-bilight graph with $n\ge3$ vertices and at least $4n-8$ edges contains $\Ke$ as a minor.
\end{theorem}

\begin{lemma}\label{lem:stability}
\begin{enumerate}
\item In a $4$-light 5-rooted graph, reducing an inclusionwise maximal reducible fragment preserves $4$-lightness.
\item Let $G$ be an unrooted $4$-bilight graph, and let $b\in\{3,4\}$. Choose $Y$ inclusionwise maximal among reducible fragments satisfying $|V(G)\setminus Y|\ge b$. Then its reducent is $4$-bilight.
\end{enumerate}
\end{lemma}

\begin{equation}\label{eq:criticaldensity}
\rho_4(G)=-9,\qquad |V(G)|\ge7.
\end{equation}

\begin{lemma}\label{lem:criticalirreducible}
The graph $G$ is $4$-irreducible.
\end{lemma}

\begin{proof}
Suppose otherwise, and choose an inclusionwise maximal reducible fragment $Y$ subject to $|V(G)\setminus Y|\ge4$. Put $Q=\dd_GY$ and let $J$ be the reducent. By Lemma~\ref{lem:stability}(2), $J$ is $4$-bilight. It is a minor of $G$ and therefore has no $\Ke$ minor.

Let $\varepsilon=0$ for a clique reduction and $\varepsilon=1$ for a dart reduction, and let $a$ be the number of edges added between surviving original vertices of $Q$. A dart introduces one new vertex and its four incident edges. A clique reduction introduces no vertex. Consequently
\[
|V(J)|=|V(G)|-|Y|+\varepsilon
\]
and
\[
e(J)=e(G)-e(G[Y])-e_G(Y,Q)+a+4\varepsilon.
\]
Subtracting four times the vertex count gives
\[
\rho_4(J)-\rho_4(G)=-\rho_4(G,Y)+a\ge0.
\]
Here both $-\rho_4(G,Y)$ and $a$ are nonnegative integers.

Since $|V(J)|\ge|V(G)\setminus Y|\ge4$, Theorem~\ref{thm:aux} applies to $J$. Its absence of a $\Ke$ minor implies $\rho_4(J)<-8$, and integrality gives $\rho_4(J)\le-9$. Combining this with~\eqref{eq:criticaldensity} and the preceding inequality yields
\begin{equation}\label{eq:zero}
\rho_4(J)=\rho_4(G)=-9,\qquad
\rho_4(G,Y)=0,\qquad a=0.
\end{equation}
The last two equalities follow because their nonnegative integer contributions sum to zero. The equality $a=0$ means that the reduction has added no edge between surviving original vertices. This is the property needed below.

In particular, the graph induced by the surviving original vertices of $J$ is exactly $G-Y$. In a clique reduction, every vertex of $J$ is such a vertex, so a $K_6$ subgraph of $J$ would already be a $K_6$ subgraph of $G$. In a dart reduction the only other vertex has degree four, whereas every vertex in a $K_6$ subgraph has at least five neighbours. It cannot belong to a 6-clique, and any 6-clique avoiding it would again be present in $G$. Therefore $J$ has no $K_6$ subgraph.

Finally, $J$ has fewer vertices than $G$. A clique reduction removes the nonempty set $Y$, while a dart reduction replaces at least two vertices by one. Hence $J$ is a smaller $4$-bilight graph on at least four vertices with $\rho_4(J)=-9$, no $\Ke$ minor and no $K_6$ subgraph. This contradicts the minimal choice of $G$.
\end{proof}

\end{document}
